\documentclass[a4paper]{amsart}

\usepackage{color}
\usepackage[colorlinks=true, citecolor=red]{hyperref}  %needs to come before amsrefs package to work with \cite{}*{}
\usepackage{amssymb,latexsym}
\usepackage[alphabetic]{amsrefs}
\usepackage[mathscr]{eucal}
\usepackage{pdfsync}
\usepackage{graphicx}
\usepackage{epstopdf}
\DeclareGraphicsRule{.tif}{png}{.png}{`convert #1 `dirname #1`/`basename #1 .tif`.png}
\usepackage[all]{xy}
\usepackage{titletoc}
%\usepackage[show]{ed}
%\usepackage{txfonts}   %  \fint  for integral average

%\usepackage{showkeys}  %options notref  notcite 

%\input{rgb.tex}

%\textcolor{DarkRed}{text in DarkRed}


\CompileMatrices

\theoremstyle{plain}
\newtheorem{theorem}{Theorem}[section]
\newtheorem{lemma}[theorem]{Lemma}
\newtheorem{proposition}[theorem]{Proposition}
\newtheorem{corollary}[theorem]{Corollary}
\newtheorem*{theorem*}{Theorem}


\theoremstyle{definition}
\newtheorem{definition}[theorem]{Definition}%[section]


\theoremstyle{remark}
\newtheorem{defn}[theorem]{Definition} 
\newtheorem{example}[theorem]{Example}
\newtheorem{remark}[theorem]{Remark}
\newtheorem{discussion}[theorem]{Discussion}

%\numberwithin{equation}{section}

%\newcommand{\abs}[1]{\lvert#1\rvert}

\DeclareMathOperator{\Lip}{Lip}
\DeclareMathOperator{\dist}{dist}
\DeclareMathOperator{\expected}{\mathbb{E}}
\DeclareMathOperator{\prb}{Prob}
\DeclareMathOperator{\spt}{spt}
\DeclareMathOperator{\lap}{\Delta}
\DeclareMathOperator{\Tr}{Tr}

\raggedbottom

\addtolength{\textwidth}{2cm}
\addtolength{\oddsidemargin}{-1cm}
\addtolength{\evensidemargin}{-1cm} 
\addtolength{\topmargin}{-1cm} 
\addtolength{\textheight}{2cm}

\newcommand{\side}[1]{\footnote{\emph{Draft Comment}:  #1}}
%\newcommand{\side}[1]{} 

\def\Xint#1{\mathchoice
   {\XXint\displaystyle\textstyle{#1}}%
   {\XXint\textstyle\scriptstyle{#1}}%
   {\XXint\scriptstyle\scriptscriptstyle{#1}}%
   {\XXint\scriptscriptstyle\scriptscriptstyle{#1}}%
   \!\int}
\def\XXint#1#2#3{{\setbox0=\hbox{$#1{#2#3}{\int}$}
     \vcenter{\hbox{$#2#3$}}\kern-.5\wd0}}
\def\ddashint{\Xint=}
\def\dashint{\Xint-}




%%% BEGIN DOCUMENT
\begin{document}

\title{Interchanging Integration and Differentiation}  
\author{John E.\ Hutchinson}
\date{\today}


\maketitle

%\tableofcontents


%%%%%%%%%%%%%%%%%%%%%%%%%%%%%%%%
%%%%%%%%%%%%%%%%%%%%%%%%%%%%%%%%
%%%%%%%%%%%%%%%%%%%%%%%%%%%%%%%%
%%%%%%%%%%%%%%%%%%%%%%%%%%%%%%%%
%\section{\textbf{Introduction}} 

This is very standard.  See, for example,~\cite{Fol}*{p 56}.


\medskip In the following we have a function $ f(x,t) $ of two variables $ x \in E \subset   \mathbb{R}^d  $ and $ t \in[a,b]$.  You may want to think of the case $ d=1 $.  Note that all integrals are with respect to the $ x $ variable.


Let  
\[
F(t) = \int_E f(x,t)\, dx.
\]
In particular we assume that $ f(x,t ) $ is integrable in  $ x $ for   each $ t\in [a,b] $.  We want to know if:
\begin{enumerate}

\item  Assuming $ f(x,t) $ is continuous in $ t $ for each $ x\in E $, does it follow that $ F(t) $ is continuous?
\item Assuming $ f(x,t) $ is differentiable in $ t $ for each $ x \in E $ does it follow that $ F'(t) $ exists and 
$ F'(t) = \int_E \frac{\partial}{\partial t} f(x,t)\, dx $?
\end{enumerate}

The answer is ``yes'', provided that the functions $ f(\cdot, t) $, respectively $ \frac{\partial}{\partial t} f(\cdot  ,t)$, are uniformly integrable over $ E $.  That is, provided they are bounded by some positive integrable function independent of $t  $.


\medskip
The proof is straightforward using 
\begin{enumerate}

\item the dominated convergence theorem, and
\item    the observation that the relevant limit  exists at $ t_0 $   if it exists    for any sequence  $ t_n \to t_0   $ and is independent of the sequence.
\end{enumerate} 
 


\begin{theorem*}
Suppose $ E \subset \mathbb{R}^d  $, $ f: E\times [a,b] \to \mathbb{R} $, and $ f(x, t) $ is integrable in $ x $ for each $ t \in [a,b] $. Let
\[
 F(t) = \int_E f(x,t)\, dx .
 \]

Suppose $  f(x,t) $ is continuous in $ t $ for each $ x\in E  $.  Suppose for all $ x\in E  $ and $ t \in [a,b] $ that $ |f(x,t)|\leq g(x) $ where  $ g : E \to \mathbb{R} $  is integrable.  Then $ F(t) $ is continuous in $ t $, i.e.
\begin{equation} \label{eq1} 
     \lim_{t\to t_0} \int_E f(x,t)\, dx = \int_E f(x,t_0)\, dx \quad \forall t_0\in [a,b] .
\end{equation} 
     
Suppose    $  \frac{\partial}{\partial t} f(x ,t)  $ exists for all $ t $ and all $ x $.  Suppose for all $ x\in E  $ and $ t \in [a,b] $ that $\left| \frac{\partial}{\partial t} f(x ,t)\right| \leq k(x) $ where  $ k : E \to \mathbb{R} $  is integrable.   Then $ F $ is differentiable and 
\begin{equation} \label{eq2} 
\left.\frac{\partial}{\partial t}\right|_{t=t_0}  \int_E f(x,t)\, dx =  \int_E \left.\frac{\partial}{\partial t}\right|_{t=t_0} f(x,t)\, dx \quad \forall t_0\in [a,b] .
\end{equation} 
\end{theorem*} 

\begin{proof} Let $ (t_n)_{n\geq 1} $ be any sequence in $ [a,b] $ such that $ t_n \to t_0 $.  Let $ f_n(x) = f(x,t_n) $ and $ f(x) = f(x,t_0) $.  Since $ f_n(x) \to f(x) $ and $ |f_n(x)| \leq g(x) $ for all $ x $, where $ \int_E g <\infty $, it follows from  the dominated convergence theorem that 
\[ 
     \lim_{t_n\to t_0} \int_E f(x,t_n)\, dx = \int_E f(x,t_0)\, dx \quad \forall t_0\in [a,b] .
\]
Since the right side is independent of the sequence $ (t_n ) $, \eqref{eq1} follows.

\medskip  Again let $ (t_n)_{n\geq 1} $ be any sequence in $ [a,b] $ such that $ t_n \to t_0 $. 
Then 
\[
h_n(x) := \frac{f(x,t_n) - f(x,t_0)}{t_n - t_0} \to  \frac{\partial f}{\partial t} (x,t_0) \quad \text{as } t_n \to t_0 .
\]
By the mean value theorem, for some $ \xi_n $ between $ t_n $ and $t_0 $,
\[
|h_n(x) | = \left| \frac{\partial f}{\partial t} (x,\xi_n) \right| \leq k(x) .
\]
It follows by the dominated convergence theorem that
\begin{align*}
\int h_n(x)\,  d x    &\to  \int_E \frac{\partial f}{\partial t} (x,t_0)\, dx \quad \text{as } t_n \to t_0,  \\
\text{i.e.  } \frac{\int_E f(x,t_n) - \int_E f(x,t_0)}{t_n - t_0} &\to \int_E \frac{\partial f}{\partial t} (x,t_0)\, dx 
                                                \quad \text{as } t_n \to t_0  .
\end{align*} 
Since this is true for any sequence $ (t_n)_{n\geq 1} $ from $[a,b]$, \eqref{eq2} follows.
\end{proof}




%%%%%%%%%%%%%%%%%%%%%%%%%%%%%%%%
%%%%%%%%%%%%%%%%%%%%%%%%%%%%%%%%
%%%%%%%%%%%%%%%%%%%%%%%%%%%%%%%%
%%%%%%%%%%%%%%%%%%%%%%%%%%%%%%%%


\begin{thebibliography}{9} 



\bib{Fol}{book}{
   author={Folland, Gerald B.},
   title={Real analysis},
%   series={Pure and Applied Mathematics (New York)},
   edition={2},
%    note={Modern techniques and their applications;
%   A Wiley-Interscience Publication},
%   publisher={John Wiley \& Sons Inc.},
%   place={New York},
   date={1999},
   pages={xvi+386},
%   isbn={0-471-31716-0},
%   review={\MR{1681462 (2000c:00001)}},
}
		

	

\end{thebibliography}

\end{document} 

